\section{KESIMPULAN}

Berdasarkan analisis hasil simulasi dan evaluasi komprehensif yang telah dilakukan, penelitian ini menghasilkan tiga kesimpulan utama sebagai berikut:
\begin{enumerate}
    \item Formulasi model Markowitz ke dalam kerangka (EPG) berhasil memetakan utilitas ekonomi menjadi sistem Hamiltonian Ising secara akurat dan konsisten. Penggunaan imbal hasil logaritmik untuk matriks kovarians terbukti efektif menjaga simetri interaksi risiko ($J_{ij}$), sementara imbal hasil sederhana mempertahankan linieritas bias ekspektasi imbal hasil ($h_i$). Implementasi suku penalti anggaran ($A$) mampu mendominasi koefisien kopling sehingga mencegah sirkuit kuantum terjebak pada portofolio yang melanggar aturan diversifikasi.
    \item Penerapan ekuilibrium Nash sebagai inisialisasi awal (\textit{warm-start}) terbukti meningkatkan efisiensi dan stabilitas konvergensi algoritma VQE secara signifikan. Injeksi solusi Nash memandu sirkuit melewati area dataran tandus yang dicirikan oleh kelenyapan varians gradien, memungkinkan algoritma mencapai keadaan dasar pada kedalaman sirkuit dangkal ($L=2$ hingga $L=3$). Hal ini memicu keruntuhan Entropi Von Neumann menuju konfigurasi probabilitas tunggal yang paling optimal, mereduksi beban komputasi secara substansial.
    \item Arsitektur HEA dengan optimasi SPSA menunjukkan performa komputasi yang efisien dalam mengeksekusi optimasi portofolio pada kedua skala sistem. Algoritma GT-VQE menghasilkan pertumbuhan modal yang superior dibandingkan VQE murni dan optimasi klasik SLSQP, dengan nilai rasio Sharpe mencapai $1{,}0185$ pada sistem $N=2$ dan $1{,}0107$ pada sistem $N=4$. Temuan ini secara empiris membuktikan bahwa integrasi teori permainan dan komputasi kuantum memberikan kontribusi nyata terhadap peningkatan kualitas keputusan alokasi aset.
\end{enumerate}

